\section{Computations}\label{sec:computations}

All computations use Magma~V2.29-4 \cite{Magma}, SageMath~10.7 with MD~Sage \cite{MDSage} (commit
\texttt{6022429}), CurveArith \cite{DerickxTerao} (commit \texttt{b7c6192}) and Zywina's package
\cite{ZywinaModular} (commit \texttt{ce31d78}); the three packages are pinned as git submodules of the
accompanying repository. The repository is organized by theorem:
\begin{center}\small
\begin{tabular}{lp{0.45\textwidth}}\toprule
folder & contents\\\midrule
\texttt{data/} & one row per curve: bounds, methods, certificates or citations\\
\texttt{certificates/A\_x0\_quintic/} & Theorem A: $36$ Magma certificates and logs\\
\texttt{certificates/B\_intermediate\_quintic/} & Theorem B and Corollary E: all $1388$ pairs $(N,\Delta)$, $N\le191$ ($1011$ strict)\\
\texttt{certificates/C\_intermediate\_gonality/} & Theorem C: models, CurveArith runs, upper-bound certificates\\
\texttt{certificates/C\_43.132.9.a.2/} & \cref{prop:43}\\
\texttt{certificates/D\_x1\_gonality/} & Theorem D: Phase~1 driver, Phase~2 program, logs\\
\texttt{certificates/good\_reduction.m} & good reduction of all canonical models used (\S\ref{sec:prelim})\\
\texttt{verify.py} & checks that every row of \texttt{data/} has a passing log or a citation\\\bottomrule
\end{tabular}
\end{center}

\subsection*{Method II for \texorpdfstring{$X_1(N)$}{X1(N)}}
Phase~1 is a short Magma script that loads CurveArith and calls its internal function
\texttt{PrecomputePowerseriesExpansions} with exactly the places and precisions that
\texttt{CAHasFunctionOfDegreeAtMost} would use; no CurveArith code is modified. It runs in several Magma processes. Every
process enumerates all places, as CurveArith does, and computes the expansions only at the places assigned to it: the
place's key (the irreducible polynomial $g(t)\in\Fp[t]$ below it, encoded as an integer, or $0$ at infinity) is
hashed by a fixed integer function, and the hash determines the process. This assignment does not depend on the order in which Magma lists the places, so
the processes together treat every place exactly once; the merge checks that the counts add up, and that the
processes agree on the curve, the place counts and the basis of differentials. (Splitting by position in the list
instead is not safe: Magma does not always list the places in the same order in different processes.) Its
output is the list of matrices of \cref{lem:rank}.

Phase~2 is a C++ program of about $1000$ lines (with GMP) that enumerates $\mathcal F_b(d)$, performs the
incremental Gaussian elimination of \S\ref{sec:phase2} over $\Fp$, and enumerates the kernel vectors up to scalars
(\cref{lem:kernel}). It is a reimplementation, for speed, of the option \texttt{Al:="Cover"} of a modified
CurveArith, which we used to cross-check it on small curves. The baseline $b$ is chosen to minimize the number of
divisors times the size of the kernels. For any $b$ the result is correct; it only affects the running time.

\subsection*{Validation}
The pipeline was tested in both directions against CurveArith's unmodified \texttt{CAGonality}: on
$X_1(26)$ over $\F_3,\F_5,\F_7$ and on $X_1(29)$ over $\F_2,\F_3$, it finds no function of degree
$\gon_{\F_p}-1$ and finds one of degree $\gon_{\F_p}$. When a function is found, its divisor is checked independently by a rank computation.
Each Phase~1 process also computes a few Riemann--Roch dimensions with Magma's \texttt{Dimension}, for
divisors supported on its own places; Phase~2 checks them against \cref{lem:rank} before searching.

\subsection*{Running times}
The $36$ certificates of Theorem~A take $40$ minutes of CPU time in total, the longest ($N=86$) $16$ minutes.
The certificates of \cref{prop:43} take a few minutes. The lower bounds of Theorem~C take about $2$ hours in
total, dominated by $48.384.19.bc.1$ ($66$ minutes). The running times for Theorem~D are in \cref{tab:D}.
